\chapter{Troubleshooting}
\label{app:trouble}

\section{Building}

\begin{description}
\item[\ctp{fixed object 0 landed on pointer 0, expected 2}] You asked
  \ct{OM=bb} to bootstrap an image. It cannot---it exists to load the 1983
  Xerox image. Build with the default \ct{OM=mt}.
\item[\ctp{SDL3: not found -- headless build}] No SDL3 development package.
  Everything works except the window. Install \ct{SDL3-devel} or
  \ct{libsdl3-dev} and rebuild.
\item[Garbled glyphs on screen] A stale object file in the archive.
  \ct|make clean && make|.
\item[The VM says the font and the image disagree] You ran \ct{make font} and
  did not rebuild the image. Some classes fix their line grid at
  class-initialisation time, so an image built against one face cannot be shown
  another. Rebuild the image.
\end{description}

\section{The window}

\begin{description}
\item[Nothing happens when I choose \ctp{browser}] It is waiting for you to
  frame the window. Press the red button, drag out a rectangle, release. See
  Chapter~\ref{ch:tenminutes}.
\item[The menu flashes and disappears] You clicked. Press and \emph{hold}.
\item[No mouse button does anything] The display was resized without the image
  being told, so no controller believes the cursor is inside it. Run with
  \ct|ST_DISPLAY_TRACE=1| to see which step refused. As a workaround,
  \ct|ST_DISPLAY_FIT=off| keeps the image's own screen size.
\item[The screen does not update when I switch back to it] Fixed; if you see
  it, the window-exposed event is not arriving. \ct{restore display} on the
  desktop menu redraws everything.
\item[The rubber band blinks while I drag] Fixed. The 1983 idiom draws the
  rectangle as a flash between two cancelling blits, and the window is now
  presented at the right moment rather than at random.
\item[The window is tiny / has huge borders] \ct|ST_DISPLAY_WINDOW=WxH| sets
  the window size, \ct|ST_DISPLAY_SCALE| the pixel scale.
\end{description}

\section{Writing code}

\begin{description}
\item[\ctp{must be a boolean}] The receiver of an inlined \ct{ifTrue:},
  \ct{and:} or \ct{whileTrue:} was not \ct{true} or \ct{false}---almost always
  an uninitialised instance variable holding \ct{nil}.
\item[\ctp{doesNotUnderstand: \#ifTrue:}] The same thing where the send was
  not inlined.
\item[A global is \ctp{nil} that should not be] Look in \ct{Undeclared}. A
  reference to a name nothing defines compiles quietly, holding \ct{nil}. A
  typo in a class name does exactly this.
\item[My \ctp{Set} lost an object I put in it] Either you overrode \ct{=}
  without \ct{hash}, or you mutated an object while it was an element.
\item[\ctp{3 + 4 * 2} answered 14] Binary selectors have no precedence among
  themselves. Parenthesise.
\item[A cascade went to the wrong receiver] A cascade sends to the receiver of
  the \emph{last message before the first semicolon}. Add \ct{yourself} if you
  want the receiver back.
\item[\ctp{10 / 4} answered \ctp{(5/2)}] \ct{/} answers an exact
  \ct{Fraction}. Use \ct{//} for integer division or \ct{asFloat}.
\item[\ctp{more than 63 literals referenced}] A method may reference at most
  63 literals; the header holds the count in six bits. Every distinct string,
  symbol, large integer and literal array counts, and repeated strings are not
  shared---two occurrences of \ct|'x'| are two mutable objects. Split the
  method, or move the data into a collection built once.
\item[\ctp{a method overflowed its frame}] The context was sized from the
  compiler's estimate of the working stack and the method pushed past it.
  Rare, and worth reporting along with the method that did it.
\end{description}

\section{Concurrency}

\begin{description}
\item[It works at one worker and fails at eight] It is a race. Read
  Chapter~\ref{ch:contract}, then look for an unsynchronised
  read-modify-write, a shared collection, or a lazily-initialised class
  variable.
\item[A process is waiting forever] If it is on a \ct{Semaphore} you are using
  for mutual exclusion, you may have taken it twice. Use \ct{Mutex}, which
  detects that and raises rather than deadlocking.
\item[The Transcript output is interleaved] Expected. Build the line in a
  \ct{WriteStream} and \ct{show:} it in one send, or hold the monitor.
\item[\ctp{Processor workerCount} answers 1] Expected: neither
  \ct{-bootstrap -eval} nor \ct{-run} starts the worker pool today. The
  parallel test suites and \ct{make bench} do.
\end{description}

\section{Saving}

\begin{description}
\item[The Browser shows no source for my methods] The image has been moved
  away from its \ct{.changes} file. Source is stored as a pointer into that
  file, not in the image.
\item[The display came back 992 wide instead of 1000] The image's own
  arithmetic: \ct{DisplayScreen} rounds a Form's width down to a sixteen-bit
  word boundary when saving shrinks and restores the display. Harmless.
\end{description}

\section{Getting more information}

\begin{center}
\small
\begin{tabular}{@{}ll@{}}
\toprule
\ctp{ST\_DISPLAY\_TRACE=1} & why a resize was refused; draw and present counts\\
\ctp{ST\_EVAL\_TRACE=1} & every bytecode an \ctp{-eval} runs\\
\ctp{ST\_BOTTOM\_LOG=1} & a return off the bottom of a stack\\
\ctp{./st80 -disasm} & what the compiler actually emitted\\
\ctp{make ASAN=1 test} & memory errors\\
\ctp{make TSAN=1 test} & data races\\
\bottomrule
\end{tabular}
\end{center}
